\begin{abstract}
Selecting a statistical procedure is scientifically valid only when the
available evidence supports its assumptions.  We study this validity-sensitive
decision using Evidence-Routed Inference (ERI), which connects diagnostics to
method eligibility or abstention.  On a frozen 18-cell Sharpe-inference pilot,
we compare four workflows that progressively vary concise versus enriched
instructions, two-call versus four-call unified execution, and a
role-specialized four-call design.  Enriched unified instructions yield the
clearest and most stable improvement over concise prompting, reducing invalid
method selection and routing regret across both scenario families and after
trimming.  Unified four-call execution has the strongest aggregate profile,
but its incremental advantage over matched two-call instructions is
concentrated in one severe structural-break cell, reverses after trimming, and
is not uniform across scenario families.  The tested specialized workflow
performs worse than the unified four-call workflow, although their bundled
differences preclude attributing this result solely to personas.  A
benchmark-specific deterministic policy provides a secondary boundary result:
when evidence and admissibility rules can be completely formalized, runtime
language-model judgment may be unnecessary.  The findings are
pilot-descriptive and motivate controlled, failure-sensitive attribution
rather than broad claims about single- or multi-agent systems.
\end{abstract}
